% =============================================================================
% Experimental Protocol: gen_3 Neural Track Extrapolation for LHCb Allen HLT1
%
% Author: G. Scriven
% Date: April 2026
%
% Scope: forward-looking protocol for the third iteration. Priority is Allen
% HLT1 deployability; Gen-3 takes the best model per family from Gen-2 and
% isolates theoretical and data-side improvements within Allen constraints.
% =============================================================================
\documentclass[a4paper,11pt]{article}
\usepackage[margin=2.5cm]{geometry}

\usepackage{amsmath,amssymb}
\usepackage{graphicx}
\usepackage{booktabs}
\usepackage{hyperref}
\usepackage{xcolor}
\usepackage{array}
\usepackage{bm}
\usepackage{pifont}
\usepackage{microtype}

% siunitx not available -- same fallback macros as gen2 protocol
\newcommand{\SI}[2]{#1~#2}
\newcommand{\SIrange}[3]{#1--#2~#3}
\newcommand{\num}[1]{#1}
\newcommand{\si}[1]{#1}
\newcommand{\kilo}{k}
\newcommand{\byte}{B}
\newcommand{\tesla}{T}
\newcommand{\mega}{M}
\newcommand{\hertz}{Hz}
\newcommand{\giga}{G}
\newcommand{\eV}{eV}
\newcommand{\mm}{mm}
\newcommand{\qop}{q/p}
\newcommand{\dz}{\Delta z}
\newcommand{\Bfield}{\mathbf{B}}
\newcommand{\state}{\mathbf{y}}
\newcommand{\Jac}{\mathbf{J}}
\newcommand{\loss}{\mathcal{L}}
\newcommand{\btheta}{\bm{\theta}}
\newcommand{\cmark}{\checkmark}
\newcommand{\xmark}{\ding{55}}

\begin{document}

\title{Experimental Protocol: Third-Generation Neural Track Extrapolation\\
for the LHCb Allen HLT1 Trigger\\
\large Allen-first iteration on Gen-2 best-of-family models}

\author{G.~Scriven\\\small Maastricht University, UHasselt, Nikhef}
\date{\today}

\begin{abstract}
This document specifies the \texttt{gen\_3} protocol for neural track
extrapolation in the LHCb experiment.  Gen-3 is a targeted iteration on
the Gen-2 best-of-family models with a single overriding design
priority: \emph{delivering an extrapolator binary that passes the Allen
HLT1 standalone smoke test}.  The Gen-2 audit
(\texttt{gen2\_new\_data\_results.tex}, §5) showed that all eleven
Gen-2 models fail Allen deployment, not because of architectural
limits, but because of three data-generation bugs (C1--C3) and two
model-level defects (A2, A4) that propagation accuracy numbers alone
cannot detect.  Gen-3 freezes three model candidates --- one per
family --- fixes the five blocking issues in dependency order, and
adds four theoretical improvements derived directly from the Gen-2
SciFi failure diagnostics.  The protocol specifies the data
regeneration, the training changes, the Allen-side validation
sequence, the hard gates each candidate must clear, and the
milestones that gate deployment.
\end{abstract}

\maketitle

% ============================================================================
\section{Introduction and priority statement}\label{gen3:sec:intro}

Gen-1 established that a $[64,64]$ MLP extrapolator can achieve
\SI{0.49}{\mm} mean position error within the \SI{64}{\kilo\byte} CUDA
\texttt{\_\_constant\_\_} memory budget.  Gen-2 extended this to five
MLP sizes, three PINN\_v2 configurations, and three NeuralRK4
configurations on a new \num{50e6}-track dataset with $\dz \in
[25,10000]$\,mm, and performed a full detector-level
characterisation against the VELO, UT and SciFi stations.

The Gen-2 conclusion is unambiguous: \textbf{the bottleneck is not the
neural architecture}.  The accuracy-ceiling Gen-2 candidate
(\texttt{nrk4\_small\_1step}, \num{17925} float32 params $=$
\SI{71.7}{\kilo\byte}) achieves \SI{0.125}{\mm} mean position error;
the deployable candidate that simultaneously satisfies the A2
($q/p$-decorrelation) and A3 (\SI{64}{\kilo\byte}) gates is
\texttt{nrk4\_tiny\_1step} (\num{4869} params $=$ \SI{19.5}{\kilo\byte},
$\langle|\Delta x|\rangle \approx \SI{0.13}{\mm}$).  What blocks
deployment is

\begin{enumerate}
  \item the training dataset uses a $qop$ convention incompatible
        with Allen (C1),
  \item the training dataset does not cover the signed-$\dz$ and
        wide input-coverage regimes that the Allen Kalman filter
        actually queries (C2, C3), and
  \item the best-accuracy candidate either overshoots the memory
        budget (A3) or carries an uncorrected deterministic
        $q/p$-linear error at the SciFi stations (A2).
\end{enumerate}

\noindent\textbf{Gen-3 priority.}  Allen HLT1 deployability is the
\emph{primary} design variable; propagation accuracy is secondary.
Every protocol decision below is justified against the Allen
\texttt{ML\_research/README.md} convention requirements and the three
smoke-test gates, not against mean position error.

\subsection{Research questions}

\noindent\textbf{RQ1 --- Deployment.}  After fixing the five blocking
issues identified in Gen-2, does at least one best-of-family model
pass the three Allen smoke-test gates
($\chi^2/\mathrm{ndof}<2.0$, $|\delta p/p|<1\,\%$,
$\sigma(\mathrm{pull}_x)<1.5$) when connected to the standalone
Kalman filter?

\noindent\textbf{RQ2 --- Which family wins in-situ.}  Which of
\{NeuralRK4, PINN\_v2, MLP\} generalises best to the in-situ Allen
distribution (signed $\dz$, $\pm3200$\,mm inputs, Allen $qop$), as
opposed to the restricted distribution the Gen-2 models were trained
on?

\noindent\textbf{RQ3 --- Theoretical improvements.}  Do the four
targeted theoretical changes (Sec.~\ref{gen3:sec:theory}) --- (i)
fixing the PINN\_v2 $q/p$-linear magnet-kick bug, (ii) adding
$z_\mathrm{start}$ to the input, (iii) detector-$\sigma$-weighted
loss, (iv) optional Jacobian regularisation against finite-difference
RK4 --- each individually reduce at least one of the Gate-1/2/3
metrics on the same model?

\section{Allen deployment constraints (hard gates)}\label{gen3:sec:allen}

All constraints in this section come from
\texttt{Allen/ML\_research/README.md} and
\texttt{Allen/ML\_research/gen2\_true\_data\_deep\_dive.md}.  They are
\emph{mandatory}; a Gen-3 candidate does not ``partially satisfy''
them.

\subsection{Conventions (C1--C3)}

\begin{description}
  \item[C1 --- $qop$ convention.] The feature-space variable the
        model consumes must be
        \begin{equation}
          qop \;=\; c_\text{light} \cdot q / p_\text{MeV},
          \qquad c_\text{light} = 299.792458\ \mathrm{mm\,ns^{-1}\,eplus^{-1}}.
          \label{eq:qop}
        \end{equation}
        This is Allen's canonical unit; a factor of $\sim300$ larger
        than the Gen-2 training convention.  The fix is in the data
        generator; no runtime conversion inside the model loader is
        permitted.
  \item[C2 --- signed $\dz$.]  The Allen Kalman filter performs
        forward \emph{and backward} passes.  The training
        distribution must therefore cover both signs:
        \begin{equation}
          \dz \sim \text{sign}\cdot\text{LogU}(25,\,10000)\ \mathrm{mm},
          \qquad \text{sign} \sim \{-1,+1\}\ \text{(balanced)}.
          \label{eq:dz_signed}
        \end{equation}
        Log-uniform in $|\dz|$ over-represents the short VELO
        inter-sensor steps (\SIrange{25}{30}{\mm}).
  \item[C3 --- input coverage.]  Inference-time inputs reach $\pm
        3200$\,mm after the first forward Kalman update; training
        must cover at least this range.  Two acceptable strategies:
        \emph{(a)} harvest $(x_\mathrm{in},y_\mathrm{in})$ pairs from
        a real Allen reconstruction at every layer interface, or
        \emph{(b)} uniform sampling over the detector envelope
        $|x|\le3500$\,mm, $|y|\le2500$\,mm.  Strategy~(a) is
        preferred because it preserves the joint density on
        $(x,y,t_x,t_y,q/p)$.
\end{description}

\subsection{Allen binary contract (A-series)}

\begin{description}
  \item[A3 --- memory budget.] Parameters + biases must fit in
        \SI{64}{\kilo\byte} CUDA \texttt{\_\_constant\_\_} memory,
        i.e.\ $\le$\num{16384} \texttt{float32} scalars.  Quantisation
        to \texttt{float16} doubles the budget but must be preceded
        by a numerical-stability check
        (Sec.~\ref{gen3:sec:validation}).
  \item[A4 --- Jacobian correctness.]  The $5\times 5$ Jacobian
        $\Jac = \partial\state_f/\partial\state_0$ with
        $\state = (x,y,t_x,t_y,qop)$ (Allen Kalman state;
        the $z_0$ and $\dz$ inputs of Fix~H are book-keeping
        variables and are \emph{not} differentiated) must agree
        element-wise with a finite-difference RK4 reference.
        Per-input FD step sizes follow Allen's own
        \texttt{compute\_jacobian} convention
        $\bm{\epsilon} = (10^{-3}\,\mathrm{mm},\,10^{-3}\,\mathrm{mm},\,10^{-6},\,10^{-6},\,\max(|qop|\cdot 10^{-4},\,10^{-8}))$.
        Structural-zero elements (where $|\Jac^\mathrm{RK4}_{ij}| <
        \epsilon^\mathrm{abs}_i$, with $\epsilon^\mathrm{abs} =
        (10^{-6},10^{-6},10^{-8},10^{-8},10^{-10})$ for the five
        output rows) are scored on absolute error $|J^{\btheta} -
        J^\mathrm{RK4}| < 3\epsilon^\mathrm{abs}_i$; the remainder
        are scored on relative error.  Median relative error must be
        $<10\,\%$ and 99th-percentile $<30\,\%$.  Frobenius-norm
        agreement is \emph{not} sufficient (Gen-1/Gen-2 lesson).
  \item[A5 --- output dimension.]  The Allen
        \texttt{MLPExtrapolator.cuh} V2 format expects \emph{five}
        outputs $(x_f,y_f,t_{x,f},t_{y,f},qop_{\!f})$.  Gen-2 models
        output four.  Gen-3 candidates must emit the fifth component;
        for NeuralRK4 this is identity ($qop$ conserved in vacuum),
        for PINN\_v2 and MLP it is an additional regression head.
  \item[A6 --- metadata block.] Each exported binary carries a
        provenance block: \texttt{qop\_convention},
        \texttt{c\_light\_value}, \texttt{dz\_signed},
        \texttt{feature\_order}, \texttt{polarity} (\texttt{MagDown}
        only in baseline Gen-3), \texttt{loader\_version\_required},
        \texttt{data\_sha256}, \texttt{git\_hash},
        \texttt{training\_seed}.  Mismatch against the Allen Kalman
        state at load time is a hard error.
  \item[A7 --- loader V3 patch.]  The existing Allen loader
        \texttt{MLPExtrapolator.cuh} hard-codes \texttt{float
        input[6]} and \texttt{float output[4]} in
        \texttt{propagate()} / \texttt{propagate\_multistep()}; it
        reads \texttt{input\_size} from the binary header and so
        does \emph{not} crash on $7$-input / $5$-output models, but
        silently reads uninitialised stack memory for the 7th input
        and silently discards the 5th output.  Gen-3 therefore
        delivers a V3 patch to \texttt{MLPExtrapolator.cuh} (inputs
        $(x,y,t_x,t_y,qop,z_0,\dz)$ in that order; outputs
        $(x_f,y_f,t_{x,f},t_{y,f},qop_f)$) \emph{as part of the
        milestone M0.5}.  A binary advertising
        \texttt{loader\_version\_required=V3} refuses to load on a
        V2-only Allen checkout.
\end{description}

\subsection{Smoke-test gates (G1--G3)}

Defined in \texttt{ML\_research/README.md}; baseline is the
production RKN extrapolator.

\begin{center}
\begin{tabular}{lll}
\toprule
Gate & Metric & Threshold \\
\midrule
G1 & $\langle\chi^2/\mathrm{ndof}\rangle$ over standalone Kalman fit & $<2.0$ \\
G2 & $|\langle p_\mathrm{fit}/p_\mathrm{true}-1\rangle|$            & $<1\,\%$ \\
G3a & $\sigma(\mathrm{pull}_x)$                                      & $<1.5$ \\
G3b & $|\langle\mathrm{pull}_x\rangle|$                             & $<0.1$ \\
\midrule
RKN baseline & ($0.95\pm0.20$,\; $-0.04\,\%$,\; $0.95$) & PASS \\
\bottomrule
\end{tabular}
\end{center}

G3 is split into G3a (width) and G3b (bias): a model with
$\langle\mathrm{pull}_x\rangle=-2$, $\sigma=1.4$ would pass the
width-only version of G3 but is clearly mis-calibrated.  G3b
catches this.

% ============================================================================
\section{Selected best-of-family candidates}\label{gen3:sec:candidates}

Gen-3 iterates on three model candidates drawn from the Gen-2
detector-level audit (\texttt{gen2\_new\_data\_results.tex}, §4.5, §5).
Rationale for each:

\begin{description}
\item[Primary: \texttt{nrk4\_tiny\_1step}.]  \num{4869} params,
      \SI{19.5}{\kilo\byte} (\emph{fits A3 with headroom}),
      $\rho(\Delta x,q/p)=+0.02$ at SciFi T3 (\emph{passes A2}),
      $\langle|\Delta x|\rangle \approx \SI{0.13}{\mm}$ one-shot.
      The \emph{only} Gen-2 model that satisfies both A2 and A3
      without further intervention; first candidate to attempt the
      smoke test.
\item[Accuracy ceiling: \texttt{nrk4\_small\_1step}.]  \num{17925}
      params, \SI{71.7}{\kilo\byte} (\emph{violates A3 by 12\,\%}),
      $\rho=+0.03$ (\emph{passes A2}), $\SI{4.9}{\mu m}$ at VELO exit,
      $\SI{25.3}{\mu m}$ at UT entry.  Deployment via
      \texttt{float16} quantisation (\SI{35.9}{\kilo\byte}); provides
      an upper bound on what the NeuralRK4 family can deliver.
\item[PINN representative: \texttt{pinn\_v2\_lam0.1}.]  \num{68612}
      params, best Gen-2 SciFi position accuracy.  Carries the
      $\rho(\Delta x,q/p)=+0.92$ magnet-kick bug
      (diagnostic ii of Gen-2 §4.5.1); blocked on A2 until the PDE
      residual is re-audited.  Kept in Gen-3 to test whether fixing
      the physics-residual formulation (Fix~G below) can rescue the
      physics-informed route.
\item[Reference only: \texttt{mlp\_medium}.]  \num{101k} params,
      best Gen-2 MLP on validation loss but fails A3 by a large
      factor.  Retained as a no-physics control; not a deployment
      target.
\end{description}

The legacy RK\_PINN, legacy PINN, and the remaining four MLP sizes are
\emph{not carried forward} --- Gen-2 already showed they are
dominated by the four above on at least one of the three Allen gates.

% ============================================================================
\section{Data improvements: fixing C1--C3 and A5}\label{gen3:sec:data}

The Gen-3 training dataset \texttt{train\_50M\_gen3.npz} is
regenerated from scratch.  Changes relative to
\texttt{train\_50M\_dz25.npz}:

\subsection{Generator patches}

\begin{enumerate}
  \item \textbf{C1 --- $qop$ convention.}  Replace the scalar
        computation in \texttt{utils/rk4\_propagator.py}:
        \begin{verbatim}
# old
qop = charge / p_MeV
# new  (Gaudi mixed units: (mm ns^-1 eplus) / MeV)
qop = 2.99792458e+2 * charge / p_MeV   # ~= 3e-2 at p = 10 GeV
        \end{verbatim}
        Sanity arithmetic: $c_\text{light}/\SI{10000}{MeV} =
        299.792458/10\,000 \approx 3.00\times 10^{-2}$, matching the
        \texttt{ML\_research/README.md} worked example.  The
        previous draft coefficient \texttt{2.99792458e-1} was a
        factor $10^{3}$ too small --- worse than the Gen-2
        convention.  \textbf{Companion edit in the integrator}:
        replace \texttt{C\_LIGHT = 2.99792458e-4} in
        \texttt{utils/magnetic\_field.py} with \texttt{C\_LIGHT =
        1.0e-6} so that $\kappa = c_\text{light}\cdot qop$ evaluates
        to the \emph{same} number as the Gen-2 expression
        $\kappa = (2.99792458\times 10^{-4})\cdot (q/p_\text{MeV})$,
        i.e.\ $c_\text{light}$ is absorbed into $qop$ exactly once.
        \textbf{Regression test (closed-form).}  Uniform
        $B_y=-1\,\mathrm{T}$, $B_x=B_z=0$, $q=+1$,
        $p=\SI{10}{GeV}$, $\dz=\SI{1000}{mm}$, initial
        $(x_0,y_0,t_{x,0},t_{y,0})=(0,0,0,0)$: the analytic
        trajectory is a circular arc of radius $R = p/(cB) =
        33.356\,\mathrm{m}$ giving $x_f = R(1-\cos(\dz/R)) =
        \SI{14.990}{\mu m}$ to six figures. Pass tolerance
        $|\Delta x_f|<\SI{1}{\mu m}$ at this test point, spot-checked
        on 10\,k tracks.  A second closed-form test with
        $B_x=+1\,\mathrm{T}$ (everything else zero) isolates the
        $B_x$ branch of the Lorentz RHS --- required because Fix~G
        depends on all three field components being correctly
        evaluated.
  \item \textbf{C2 --- signed $\dz$.}  Replace
        $\dz \sim \mathcal{U}(25,10000)$ with
        \begin{equation}
          \dz = \epsilon \cdot 10^{u},
          \quad \epsilon\sim\{-1,+1\},\ u\sim\mathcal{U}(\log_{10}25,\,\log_{10}10000).
        \end{equation}
        The RK4 integrator handles negative $\dz$ by stepping with
        $-h$; the magnetic-field module is independent of step sign.
  \item \textbf{C3 --- full input coverage.}  Extend the sampling
        envelope to
        $x_0\sim\mathcal{U}(-3500,+3500)$\,mm,
        $y_0\sim\mathcal{U}(-2500,+2500)$\,mm, and widen the slopes
        to $t_{x,0}\sim\mathcal{U}(-0.4,0.4)$,
        $t_{y,0}\sim\mathcal{U}(-0.35,0.35)$ to accommodate Kalman
        states that extrapolate \emph{off-axis and wide-angle}
        simultaneously.  Retain the $|x|,|y|<\SI{5000}{\mm}$,
        $|t|<2.0$ post-integration rejection.  The log-momentum
        range $p\in[1,200]$\,GeV/$c$ is unchanged.
  \item \textbf{A5 --- fifth output.}  Add $qop_f$ to the output
        array $\mathbf{Y}\in\mathbb{R}^{N\times 5}$.  For vacuum
        propagation $qop_f\equiv qop_0$; this column is kept
        explicitly so the Allen binary contract is honoured without
        a loader-side synthetic column.
\end{enumerate}

\subsection{Data provenance hardening}

Alongside the regenerated \texttt{.npz} file, write
\texttt{train\_50M\_gen3.meta.json} capturing
\texttt{qop\_convention=allen\_v1},
\texttt{dz\_signed=True},
\texttt{x\_range=[-3500,3500]},
\texttt{generator\_git\_hash},
\texttt{field\_map\_sha256},
\texttt{batch\_seeds=\{42 + 7919 i\}},
\texttt{n\_batches},
\texttt{created\_utc}.  This metadata is passed through to the Allen
binary's A6 block on export.

\subsection{Optional: realism extensions (deferred to Gen-3.1)}

The following are \emph{not} in the baseline Gen-3 dataset but are
staged for Gen-3.1 once the Allen smoke test passes on the vacuum
dataset:

\begin{itemize}
  \item Multiple scattering: Highland formula applied as a random
        slope kick at each RK4 sub-step.
  \item Energy loss: Bethe--Bloch $dE/dx$ applied between sub-steps;
        makes $qop_f \ne qop_0$ and activates the A5 output.
  \item Mixed magnet polarity (MagUp via $B_y \to -B_y$).
\end{itemize}

% ============================================================================
\section{Theoretical improvements}\label{gen3:sec:theory}

Each item is derived from a specific diagnostic in Gen-2 §4.5.1 and
§5.  Numbering continues from the Gen-2 fix labels (A--F).

\subsection{Fix G --- PINN\_v2 PDE residual re-audit (A2)}

Gen-2 measured $\rho(\Delta x, qop) = +0.92,\; +0.92,\; +0.99$ for
PINN\_v2 with $\lambda\in\{0.1,1.0,10.0\}$.  The correlation
\emph{grows with $\lambda$}, which rules out a pure data-fit
explanation and implicates the PDE residual itself.  Two candidate
causes:
\begin{enumerate}
  \item The field-interpolation module returns only $B_y$, and the
        Lorentz RHS (Eqs.\,\ref{eq:dtx}--\ref{eq:dty} of Gen-2
        protocol) is evaluated with $B_x = B_z = 0$ by default.  The
        measured field map has $|B_x|$ up to \SI{0.15}{\tesla} in
        the fringe regions that the SciFi-station tracks traverse.
  \item The non-dimensionalisation factor $\sigma_k/\dz$
        (Eq.~5.9 of Gen-2 protocol) does not remove the $q/p$
        dependence of the residual scale, so the penalty is
        effectively weighted by $|q/p|$.
\end{enumerate}
Action: (i) restore all three field components in
\texttt{models/pinn\_v2.compute\_physics\_loss}; (ii) unit-test the
PDE residual against the analytic uniform-$B_y$ trajectory
(closed-form circular arc) to confirm the residual is identically
zero on the analytic ground truth.  Pass criterion: $\rho(\Delta x,
qop)<0.1$ at SciFi T3 on the Gen-3 test split.

\subsection{Fix H --- add $z_\mathrm{start}$ to the input}

The field map is strongly $z$-dependent ($|B_y|$ peaks at
$z\simeq\SI{5007}{\mm}$) but the Gen-2 input vector
$\mathbf{x}=(x_0,y_0,t_{x,0},t_{y,0},qop,\dz)$ does \emph{not}
carry $z_0$.  All Gen-2 models implicitly assume $z_0=0$; the RK4
ground truth does not.  This produces a deterministic
$z_0$-dependent residual visible in the per-station error growth
(Gen-2 Tab.~\ref{gen2new:tab:error_legs}).  Gen-3 extends the input
to $\mathbf{x}\in\mathbb{R}^{7}$:
\begin{equation}
  \mathbf{x} = (x_0,\,y_0,\,t_{x,0},\,t_{y,0},\,qop,\,z_0,\,\dz).
\end{equation}
The $z_0$ column already exists in the Gen-2 $\mathbf{Z}$ array; the
regenerated Gen-3 dataset includes it in $\mathbf{X}$ directly.  Cost:
one extra first-layer weight column per neuron; negligible
memory impact.

\subsection{Fix I --- detector-$\sigma$-weighted endpoint loss}

Gen-2 uses a plain component-weighted MSE
(Eq.~5.7 of Gen-2 protocol) with weights $1/\sigma_k$ taken from the
\emph{training-set} standard deviation.  This produces uniform
tolerance across all $\dz$ and all momenta.  But the downstream
Kalman filter evaluates residuals in \emph{detector-resolution
units}:
$(\sigma^\mathrm{VELO}_x,\sigma^\mathrm{VELO}_y)=(\SI{12}{\mu m},\SI{12}{\mu m})$,
$(\sigma^\mathrm{UT}_x,\sigma^\mathrm{UT}_y)=(\SI{50}{\mu m},\SI{5}{\mu m})$ (stereo UT layers; $y$-strips),
$(\sigma^\mathrm{SciFi}_x,\sigma^\mathrm{SciFi}_y)=(\SI{60}{\mu m},\SI{500}{\mu m})$ (stereo SciFi layers).
Slope tolerances are set by differentiating position resolution over
typical inter-station lever arms: $\sigma_{t_x}=\SI{1e-4}{}$ (VELO/UT), $\SI{3e-5}{}$ (UT$\to$SciFi), $\sigma_{t_y}$ identical scale.
Gen-3 replaces the training loss with
\begin{equation}
  \loss_\mathrm{det}
  \;=\;
  \frac{1}{N}\sum_i \sum_{k\in\{x,y,t_x,t_y\}}
      \Big(\frac{\hat{y}_{ik} - y_{ik}}{\sigma^\mathrm{det}_k(z_f^i)}\Big)^{\!2},
  \label{eq:det_loss}
\end{equation}
where $\sigma^\mathrm{det}_k(z)$ is the detector resolution of the
\emph{nearest station} to $z_f$.  This directly couples the training
objective to the quantity Allen Gate~G3 measures.  Expected effect:
pull-width narrowing at VELO and UT at the cost of mild degradation
in the SciFi-interior region where no station provides a hit.

\subsection{Fix J --- Jacobian regularisation via finite-diff RK4
(A4)}

Gen-2 §6 (Jacobian Frobenius norms) showed all models sit within a
factor of~2 of the RK4 reference on $\|\Jac\|_F$, but this is a
scalar summary; element-wise Jacobian agreement was never measured.
The Kalman covariance update is linear in the Jacobian, so an
element-wise bias of even \num{20}\,\% can inflate $\sigma(\mathrm{pull}_x)$
past Gate~G3.  Gen-3 adds an optional regulariser:
\begin{equation}
  \loss_\mathrm{Jac}
  \;=\;
  \lambda_J \cdot
  \bigg\|
     \frac{\partial f_{\btheta}}{\partial \mathbf{x}_0}
     \,-\,
     \frac{\partial f_\mathrm{RK4}}{\partial \mathbf{x}_0}
  \bigg\|_F^{\,2},
  \label{eq:jac_loss}
\end{equation}
with $\partial f_{\btheta}/\partial\mathbf{x}_0$ computed by
\texttt{torch.func.jacfwd} and the RK4 Jacobian by centred finite
differences using Allen's per-input step sizes
$\bm{\epsilon}=(10^{-3},10^{-3},10^{-6},10^{-6},\max(|qop|\cdot 10^{-4},10^{-8}))$
(\emph{not} a flat $10^{-4}$, which would be noise-dominated on the
slope and $qop$ rows).  Activated on a \num{5e3}-sample mini-batch every
\num{10} optimiser steps to amortise the finite-difference cost.
$\lambda_J$ swept in $\{0,\,10^{-3},\,10^{-2},\,10^{-1}\}$.
The 10-step amortisation cadence is a throughput/signal trade-off:
every optimiser step logs $\loss_\mathrm{Jac}$ (even on steps where
it is not applied) via a cached previous-value so that post-hoc
diagnostic plots can detect drift between the data gradient and the
Jacobian gradient.

\subsection{Fix K --- optional: predict residual to straight-line
ansatz}

For the MLP family specifically, the PINN\_v2 envelope trick
(prediction = straight line + learnable correction) provides a
\num{10}\,\%-level accuracy improvement when the straight-line
baseline already captures the ballistic component.  Gen-3 adds an
ablation in which \texttt{mlp\_medium} is trained to emit
$\bm{\delta}=\hat{\state}-\state_\text{base}$ instead of
$\hat{\state}$ directly.  This is a pure output-space transformation;
no architectural change.

\subsection{Fix L --- optional: non-uniform multi-step NeuralRK4}

\texttt{nrk4\_small\_2step} underperformed \texttt{1step} on SciFi
($\rho=-0.49$); the suspected cause is that the step split is
uniform in $\dz$, so both RK4 sub-steps see similar field gradients.
A non-uniform split with an internal breakpoint at the magnet peak
$z\simeq\SI{5000}{\mm}$ is staged as a Gen-3.1 ablation if
\texttt{1step} fails at the SciFi-interior.

\section{Model-level changes per candidate}\label{gen3:sec:model_changes}

\begin{center}
\small
\setlength{\tabcolsep}{4pt}
\begin{tabular}{l p{4.3cm} p{4.3cm} p{3.2cm}}
\toprule
Model & Gen-3 architectural change & Gen-3 loss change & New params \\
\midrule
\texttt{nrk4\_tiny\_1step} &
  $\mathbf{x}\in\mathbb{R}^{7}$ (Fix H); 5-output head (A5). &
  Fix I; Fix J at $\lambda_J\in\{0,10^{-3},10^{-2}\}$. &
  $\approx$\num{5000}; \SI{20}{\kilo\byte}. \\
\texttt{nrk4\_small\_1step} &
  As above + float16 export. &
  Fix I; Fix J. &
  $\approx$\num{18}k; \SI{36}{\kilo\byte} (fp16). \\
\texttt{pinn\_v2\_lam0.1} &
  Fix G (re-audited PDE residual + full $\Bfield$); Fix H; 5-output head. &
  Fix I; Fix G (unit-tested residual); $\lambda_\mathrm{pde}$ swept
  $\{0.03, 0.1, 0.3\}$ only. &
  $\approx$\num{69}k; \emph{reference only}. \\
\texttt{mlp\_medium} &
  Fix H; 5-output head; Fix K ablation. &
  Fix I. &
  $\approx$\num{101}k; \emph{reference only}. \\
\bottomrule
\end{tabular}
\end{center}

% ============================================================================
\section{Training protocol}\label{gen3:sec:training}

Largely inherited from Gen-2 §6.  Explicit Gen-3 deltas:

\begin{itemize}
  \item \textbf{Split.}  80/10/10 on the regenerated dataset, seed
        42, stratified by $\text{sign}(\dz)$ so that the backward
        pass is represented in validation and test.
  \item \textbf{Optimiser.}  AdamW, $\beta=(0.9,0.999)$, weight
        decay $10^{-4}$.  Base learning rate $10^{-3}$ for
        NeuralRK4/MLP, $5\times 10^{-4}$ for PINN\_v2.  Cosine
        annealing, 5-epoch linear warm-up, max 200 epochs, early
        stop on validation loss with patience 30 (NeuralRK4/MLP) /
        40 (PINN\_v2).
  \item \textbf{Batch size.}  2048 for NeuralRK4/MLP, 1024 for
        PINN\_v2, reduced to 512 on mini-batches that evaluate
        Fix~J.
  \item \textbf{Seeds.}  Global seed 42; dataset seed 42; per-batch
        seeds $s_i = 42 + 7919 i$ (unchanged from Gen-2).
  \item \textbf{Infrastructure.}  Nikhef HTCondor cluster;
        \texttt{condor/submit\_gen3.sub}.  MLflow experiment
        \texttt{gen\_3\_track\_extrapolation}.
\end{itemize}

% ============================================================================
\section{Validation protocol}\label{gen3:sec:validation}

Gen-3 validation is a three-stage funnel.  A candidate that fails any
earlier stage does not proceed to the next.

\subsection{Stage 1 --- propagation-level sanity (reuse Gen-2
characterisation)}

On the Gen-3 test split, compute the Gen-2 detector-level metrics:
per-station $\langle|\Delta x|\rangle$,
$\rho(\Delta x,qop)$ at SciFi T3, per-leg residual growth,
two-dimensional residual heatmap in $(\dz, p)$.  Pass criteria:
\begin{itemize}
  \item VELO exit: $\langle|\Delta x|\rangle < \SI{12}{\mu m}$
        (VELO $\sigma$).
  \item UT entry:  $\langle|\Delta x|\rangle < \SI{50}{\mu m}$
        (UT $\sigma$).
  \item SciFi T3: $|\rho(\Delta x,qop)|<0.1$ for PINN\_v2; $<0.3$
        otherwise.
  \item Backward pass ($\dz<0$) metrics within a factor~2 of forward.
\end{itemize}

\subsection{Stage 2 --- Jacobian correctness (A4)}

On \num{1000} randomly selected test samples, compute the
$5\times 5$ Jacobian $\Jac_{\btheta}$ via
\texttt{torch.func.jacfwd} over the Allen Kalman state
$(x,y,t_x,t_y,qop)$ only (not $z_0$, $\dz$) and
$\Jac_\mathrm{RK4}$ by centred finite differences with Allen's
per-input step sizes (A4).  Pass criterion per A4: median
element-wise relative error $<10\,\%$, 99th-percentile $<30\,\%$
on the non-structural-zero elements; structural zeros scored on
absolute error.  Report $(\Jac_{\btheta} -
\Jac_\mathrm{RK4})/\Jac_\mathrm{RK4}$ per matrix element so that
any systematic bias is visible.

\subsection{Stage 2.5 --- exported binary round-trip}

Load the exported \texttt{.bin} into a reference
\texttt{numpy}-only loader
(\texttt{scripts/load\_allen\_v3\_ref.py}) and evaluate 100 test
tracks.  Assert bitwise agreement ($|\Delta| < 10^{-6}$ relative)
with the PyTorch forward pass on the same inputs.  This catches the
class of failure where Stage 3 regresses for loader reasons rather
than model reasons --- the single most common Gen-1/Gen-2 deployment
bug.

\subsection{Stage 3 --- Allen standalone smoke test (G1--G3)}

Export the model to the \texttt{MLPExtrapolator.cuh} V2 binary via
\texttt{scripts/export\_allen\_v2.py} and run
\texttt{ML\_research/standalone/run\_extrapolators} against the RKN
baseline.  Record G1/G2/G3 on the same sample of
$\mathcal{O}(10^5)$ tracks.  A candidate is declared \emph{deployment
ready} if and only if all three gates pass simultaneously.

\subsection{Validation harness outputs}

Each candidate produces:
\texttt{gen3\_candidate\_\{name\}\_stage1.csv} (per-station metrics),
\texttt{\dots\_stage2.csv} (Jacobian element-wise stats),
\texttt{\dots\_stage3\_smoketest.log} (Allen harness stdout),
\texttt{\dots.bin} (exported weights), and
\texttt{\dots.meta.json} (A6 provenance block).

% ============================================================================
\section{Iteration plan}\label{gen3:sec:iteration}

Dependency-ordered; later milestones assume earlier ones pass.

\begin{enumerate}
  \item \textbf{M0 --- data regeneration.} Patch generator for
        C1/C2/C3/A5, regenerate \num{50e6} tracks
        (\texttt{condor/submit\_gen3\_datagen.sub}); verify
        closed-form uniform-$B_y$ and uniform-$B_x$ tests
        (§\ref{gen3:sec:data}) with tolerance
        $|\Delta x_f|<\SI{1}{\mu m}$ at $(p,\dz)=(\SI{10}{GeV},\SI{1}{m})$.
        \emph{Gate: dataset SHA-256 recorded; both closed-form tests
        pass; spot-check 10k tracks match the Allen-convention
        reference.}
  \item \textbf{M0.5 --- Allen loader V3 patch (A7).}  Patch
        \texttt{MLPExtrapolator.cuh} \texttt{propagate()} /
        \texttt{propagate\_multistep()} to accept the
        $7$-input / $5$-output format; add the
        \texttt{loader\_version\_required} header field; add
        \texttt{compute\_jacobian} extension to the $5\times 5$
        Kalman-state Jacobian.  \emph{Gate: unit test loading a
        Gen-3 placeholder binary into a local Allen build succeeds;
        reference Python loader
        (\texttt{scripts/load\_allen\_v3\_ref.py}) and C++ loader
        produce bitwise-identical outputs on 100 test tracks.}
  \item \textbf{M1 --- baseline re-training.}  Train
        \texttt{nrk4\_tiny\_1step} with Fix~H only
        (see M1 post-mortem below; Fix~I deferred to M2, Fix~J and
        Fix~G not active).
        \emph{Gate: Stage~1 passes.}
  \item \textbf{M2 --- Jacobian validation.}  Run Stage~2 on the M1
        model.  If median element-wise error $>10\,\%$, add Fix~J
        (Jacobian regulariser) and retrain.  \emph{Gate: Stage~2
        passes.}
  \item \textbf{M3 --- first smoke test.}  Export M2 model and run
        Stage~3.  \emph{Gate: at least one of G1/G2/G3 is closer to
        baseline than the Gen-2 true-data result
        (Tab.~\ref{gen2new:tab:allen_prev} in Gen-2 report).}
  \item \textbf{M4 --- smoke-test pass.}  Iterate on $\lambda_J$ and
        hidden-dim choice until all three gates pass.  \emph{Gate:
        G1$<2.0$, $|$G2$|<1\,\%$, G3$<1.5$ simultaneously.}
  \item \textbf{M5 --- accuracy ceiling.}  Repeat M1--M4 for
        \texttt{nrk4\_small\_1step} with float16 export.  \emph{Gate:
        same three gates pass at higher accuracy margin.}
  \item \textbf{M6 --- PINN\_v2 rescue.}  Train
        \texttt{pinn\_v2\_lam0.1} with Fix~G + Fix~H + full
        $\Bfield$.  \emph{Gate: $\rho(\Delta x,qop)<0.1$; if not,
        PINN\_v2 is retired.}
  \item \textbf{M7 --- Kalman integration.}  Replace the RKN
        extrapolator in the Allen HLT1 sequence with the M4/M5
        model; record event-throughput impact and end-to-end
        physics performance (tracking efficiency, fake rate,
        momentum resolution).  \emph{Gate: no regression in
        tracking efficiency; per-track inference time $\le$ RKN
        baseline on the target HLT1 GPU (A5000 assumed; re-state
        on the actual deployment GPU if different); throughput
        improvement measurable at the event level.}
\end{enumerate}

% ============================================================================
\subsection{M1 post-mortem and gen-2 inheritance audit}
\label{gen3:sec:m1_postmortem}

An initial Gen-3 M1 attempt (commit prior to the present revision)
failed to learn: after 28~epochs on \num{1.4e6} training tracks the
\texttt{nrk4\_tiny\_1step} correction scale
($\exp \mathtt{log\_corr\_scale}$) regressed toward zero, i.e.\ the
network actively suppressed its own correction branch and fell back to
the pure RK4 baseline ($\sim\SI{164}{\mu m}$ mean $|\Delta x|$ on
signed-$\dz$ tracks).  Root-cause analysis against the Gen-2 training
recipe (\texttt{Quantum\_Track\_Reconstruction/Recreating\_paper\_results/gen2\_new\_data\_results.tex})
identified three unjustified departures from the recipe that had
shipped 0.125\,mm in Gen-2:

\begin{enumerate}
  \item \textbf{Loss change.}  Fix~I (detector-$\sigma$ loss) was
        active on M1.  Detector resolutions
        $\sigma \sim \mathcal{O}(\SI{10}{\mu m})$ at VELO make those
        residuals dominate the objective, and the pure RK4 baseline
        for large signed $|\dz|$ sits at $\sim\SI{164}{\mu m}$
        --- more than an order of magnitude above the VELO target.
        The optimiser's fastest route to loss reduction is to
        \emph{shrink} the correction branch before the network has
        learnt a useful correction; that is the observed failure
        mode.  Gen-2 used a \emph{data-driven} normalised MSE
        ($\sum_k ((y_k-\hat y_k)/\sigma^{\mathrm{data}}_k)^2$)
        which weights each output by its empirical spread and is
        therefore soft on the regions where the pure-RK4 floor is
        unavoidable.  \textbf{Resolution:} M1 reverts to the Gen-2
        normalised MSE.  Fix~I is deferred to M2, where it may be
        introduced as a secondary objective after the correction
        branch has converged.
  \item \textbf{Hardcoded $z$-normalisation.}  The NeuralRK4
        correction network previously assumed $z_0\in[0,\SI{10}{m}]$
        and used a fixed $(z-\num{5000})/\num{5000}$ normaliser.
        With the Gen-3 signed-$\dz$ data (Fix~C2) $z_0$ can be up
        to $\SI{14}{m}$, violating that assumption.
        \textbf{Resolution:} the normaliser now reads the dataset
        mean/std from the model's \texttt{input\_mean}/
        \texttt{input\_std} buffers, consistent with the treatment
        of every other input component.
  \item \textbf{Correction network blind to $\dz$.}  In Gen-2 the
        $\dz$ values were unsigned and bounded ($\in[25,
        \SI{10000}]\,\SI{}{mm}$), so the correction network's
        inputs $(x,y,t_x,t_y,z,q/p)$ were sufficient.  In Gen-3
        with signed $\dz$ the same state maps to two qualitatively
        different propagations (forward vs backward), and the pure
        RK4 residual grows as $\mathcal{O}(\dz^5)$ --- so the
        correction network \emph{must} see $\dz$ to scale its own
        output.  \textbf{Resolution:} the correction MLP input is
        extended from 6 to 8 features, adding
        $(\log_{10}|\dz|/100,\,\operatorname{sign}\dz)$.  This is
        the only new \emph{architectural} change in Gen-3 relative
        to the Gen-2 recipe; all other differences are data or
        loss.  It is covered by Fix~C2 rather than a new Fix letter.
\end{enumerate}

\noindent After these three corrections, M1 smoke tests
(\SI{80}{kilo} tracks $\times$ 3 epochs, CPU) produce:

\begin{center}
\small
\begin{tabular}{lccc}
\toprule
Model & initial pos-rms & epoch-3 pos-rms & epoch-3 $|\Delta x|$ \\
\midrule
\texttt{mlp\_small}             & \SI{2569}{mm}  & \SI{1755}{mm}   & --- \\
\texttt{nrk4\_small\_1step}     & \SI{785}{\mu m}& \SI{785}{\mu m} & \SI{47}{\mu m} \\
\texttt{pinn\_v2\_small}        & \SI{860}{\mu m}& \SI{1515}{\mu m}& \SI{615}{\mu m} \\
\bottomrule
\end{tabular}
\end{center}

\noindent The numbers confirm that (i) NeuralRK4 initialises at the
pure-RK4 floor (correct behaviour from the 1-step integrator;
correction branch gated by \texttt{log\_corr\_scale$\to\!-\infty$}
until data loss warrants activation); (ii) the MLP learns from
scratch on a non-trivial trajectory and is therefore still far from
the floor at 3 epochs, as expected; and (iii) PINN\_v2 is still
inside its physics-loss warm-up ramp ($\lambda_\mathrm{pde}$ at
0.4--0.6 of target) and shows the classic PINN early-training
divergence --- Gen-2 observed the same behaviour and recovered by
epoch $\sim 20$.  None of these smoke-test points are gated
results; full M1--M4 training is required before any verdict.

\subsubsection{M1 full-run outcome (\texttt{*\_v1}, 200\,k tracks
$\times$ 40 epochs, CPU; training 2026-04-24, audit \& writeup
2026-05-06)}
\label{gen3:sec:m1_v1_outcome}

The full M1 campaign completed on 2026-04-24.  Headline numbers
on the held-out test split:

\begin{center}
\small
\begin{tabular}{lcccc}
\toprule
Run & trainable & best ep & mean $|\Delta x|$ & wall \\
\midrule
\texttt{mlp\_small\_v1}        & 18\,076 (\SI{70.6}{\kilo\byte}) & 40/40 & \SI{18.41}{\mm}    & 6.3\,m \\
\texttt{nrk4\_tiny\_1step\_v1} &  4\,997 (\SI{19.5}{\kilo\byte}) &  4/21 & \SI{0.113}{\mm}     & 20.4\,m \\
\texttt{nrk4\_small\_1step\_v1}& 18\,181 (\SI{71.0}{\kilo\byte}) &  4/21 & \SI{0.210}{\mm}     & 20.6\,m \\
\texttt{pinn\_v2\_small\_v1}   & 10\,372 (\SI{40.5}{\kilo\byte}) & 16/25 & \SI{0.117}{\mm}     & 54.1\,m \\
\bottomrule
\end{tabular}
\end{center}

\noindent The triage of these results identified three issues
that together explain why none of them passes the Stage-1 gate:

\paragraph{F1 --- checkpoint $\neq$ weight memory (resolved, no
bug).}  All four \texttt{best\_model.pt} files are 11\,MB on disk
because the LHCb field grid (\texttt{field.B\_grid}, shape
$1\times 3\times 146\times 81\times 81$, $\sim$\SI{11}{\mega\byte})
is registered as a non-trainable buffer.  Trainable-parameter
counts are the protocol-budget numbers in the table above; the
Allen export pipeline strips the \texttt{field.*} keys so this is
not a deployment-budget concern.

\paragraph{F2 --- the data-MSE loss is heavy-tail dominated.}
Re-running \texttt{validate()} on \texttt{nrk4\_tiny\_1step\_v1}
across all three splits with the saved checkpoint gives:

\begin{center}
\small
\begin{tabular}{lccc}
\toprule
Split & loss & mean $|\Delta x|$ & RMS pos\\
\midrule
val   & $1.43\!\times\!10^{-5}$ & \SI{0.122}{\mm} & \SI{1.15}{\mm}\\
test  & $1.92\!\times\!10^{-7}$ & \SI{0.113}{\mm} & \SI{0.66}{\mm}\\
train & $2.16\!\times\!10^{-4}$ & \SI{0.125}{\mm} & \SI{2.64}{\mm}\\
\bottomrule
\end{tabular}
\end{center}

The loss varies by a factor of \num{1100} between splits; mean
$|\Delta x|$ varies by $<10\,\%$.  The cause is a heavy tail in the
residual distribution: a sub-percent fraction of high-$|q/p|$,
large-$|\dz|$, backward-extrapolated tracks produce
$\mathcal{O}(\mathrm{cm})$ errors, and the squaring in the
$\sigma_\mathrm{data}$-normalised MSE makes a single such event
worth $\sim 10^4$ typical events.  The number of tail events
varies between splits at a level that swamps the body
distribution.  \emph{Implication}: the squared-loss is unsuitable
as a model-selection metric on Gen-3 signed-$\dz$ data, and the
``best epoch by val\_loss'' criterion that worked on Gen-2 is no
longer reliable.  This is a genuine Gen-3-specific finding ---
Gen-2's bounded unsigned $|\dz|$ regime did not encounter it.
The M1.5 fix (Fix~M, below) is to switch the selection metric to
median $|\Delta x|$ and to wrap the residual in a $\log\cosh$
non-linearity that has $L^2$ behaviour on the body and $L^1$
asymptote on the tail.

\paragraph{F3 --- the correction net destabilises pure RK4 under
LR warm-up.}  The full per-epoch trace of
\texttt{nrk4\_tiny\_1step\_v1} reads (on val):

\begin{center}
\small
\begin{tabular}{ccccc}
\toprule
ep & lr & val\_loss & pos\_rms & $|\Delta x|$\\
\midrule
 1 & $1.0\!\times\!10^{-4}$ & $1.0\!\times\!10^{-5}$ & \SI{1142}{\mu m} & \textbf{\SI{47.4}{\mu m}}\\
 2 & $2.0\!\times\!10^{-4}$ & $1.0\!\times\!10^{-5}$ & \SI{1142}{\mu m} & \SI{52.9}{\mu m}\\
 3 & $3.0\!\times\!10^{-4}$ & $1.0\!\times\!10^{-5}$ & \SI{1143}{\mu m} & \SI{72.1}{\mu m}\\
 4 & $4.0\!\times\!10^{-4}$ & $1.0\!\times\!10^{-5}$ & \SI{1146}{\mu m} & \SI{121.7}{\mu m}\\
 5 & $5.0\!\times\!10^{-4}$ & $1.0\!\times\!10^{-5}$ & \SI{1174}{\mu m} & \SI{254.4}{\mu m}\\
 \multicolumn{5}{c}{\dots monotone degradation \dots}\\
21 & $2.8\!\times\!10^{-4}$ & $1.7\!\times\!10^{-5}$ & \SI{8561}{\mu m} & \SI{7.1}{\mm}\\
\bottomrule
\end{tabular}
\end{center}

\noindent Pure RK4 (epoch 1, LR still in warm-up at $10^{-4}$)
delivers \textbf{\SI{47}{\mu m}} mean $|\Delta x|$ --- already
\emph{below} the gen-2 converged \SI{125}{\mu m} reference, with
zero learnable parameters.  As the cosine schedule ramps the LR
through $5\!\times\!10^{-4}$ the correction branch fires and the
network monotonically erases this gain, ending at \SI{7.1}{\mm}
when early-stopping triggers.  The recorded ``best ep~4'' is an
artefact of F2 (val\_loss is constant in the noise floor while
mean $|\Delta x|$ is already drifting at epoch 2).

\noindent The interpretation is direct: the \emph{dz-aware}
correction-net features
$(\log_{10}|\dz|/100,\,\operatorname{sign}\dz)$ added in Gen-3
(Fix~K) span three orders of magnitude in scale, and at peak LR
the network learns large corrections appropriate for the heavy
tail and applies them to the body of the distribution.  Gen-2 with
six bounded inputs did not encounter this regime.

\subsection{M1.5 plan: principled improvement, not redesign}
\label{gen3:sec:m1_5_plan}

The M1 finding is that pure RK4 on the 7-input formulation
(Fix~H) is already \emph{better} than the gen-2 corrected
benchmark on the body of the distribution; the open problem is
the long-$|\dz|$ tail.  M1.5 is therefore aimed at \emph{not
breaking} pure RK4 on the body while letting the network correct
the tail.  The four interventions below are dependency-ordered
and individually motivated by an M1 finding.

\begin{description}
  \item[Fix~M --- $\log\cosh$ on detector-$\sigma$-normalised
    residuals + median early-stop.]
    Replace $((y-\hat y)/\sigma_\mathrm{data})^2$ with
    $\log\cosh\!\big((y-\hat y)/\sigma_\mathrm{target}\big)$, which
    is $L^2$ for small residuals (preserves the gen-2 convergence
    on the body) and $L^1$-asymptotic for large residuals (caps
    the tail's contribution to gradients).  Replace the val\_loss
    early-stop / best-epoch criterion with median $|\Delta x|$ on
    the val split.  Addresses F2.
  \item[Fix~N --- truncation-error-gated correction.]
    Replace the additive correction
    $\Delta\mathbf{x}=s\cdot \mathrm{NN}(\dots)$ with the
    multiplicative form
    $\Delta\mathbf{x}=s\cdot \mathrm{NN}(\dots)\cdot\hat\sigma_{RK4}$,
    where $\hat\sigma_{RK4}$ is an analytical local-truncation
    estimate (proportional to $|\dz|^5$ and to the FSAL
    difference $\|\mathbf{k}_4-\mathbf{k}_3\|$).  This gates the
    correction off where pure RK4 is already accurate (the body)
    and lets it act where RK4 is provably weak (the tail).
    Addresses F3 structurally.
  \item[Fix~O --- correction-input rescaling and \texttt{sign}-removal.]
    Drop $\operatorname{sign}\dz$ from the correction-net input
    (it is already encoded in the RK4 step direction).  Z-score
    the $\log_{10}|\dz|/100$ feature using the training
    statistics so the correction MLP's first-layer gradients are
    isotropic.  Reduces input-feature dimension to seven and
    removes a known source of correction-branch overshoot.
  \item[Fix~P --- separate-LR optimiser groups + grad-norm guard.]
    Optimise the correction MLP at $\mathrm{lr}_\mathrm{corr}=
    \mathrm{lr}_\mathrm{main}/5$.  Add a moving-median gradient-norm
    backstop: if $\|\nabla(\mathrm{correction\_net})\|>
    10\times\mathrm{median}$ for three consecutive batches,
    freeze \texttt{log\_corr\_scale} for one epoch.  Operational
    safety net against the F3 runaway.
\end{description}

The MLP and PINN tracks are advanced in parallel: the MLP is
deepened to a 4-layer recipe (matching gen-2 \texttt{mlp\_medium});
the PINN is held at M1 results until Fix~G's PDE-residual
unit-test (\texttt{tests/test\_pinn\_residual.py}) is implemented
and passes.  These are necessary preconditions, not part of the
M1.5 NRK4 thrust.

\subsection*{Gen-3 thesis-defensible findings (settled by M1)}

\begin{enumerate}
  \item Pure RK4 on the 7-input (Fix~H) formulation achieves
        \SI{47}{\mu m} mean $|\Delta x|$ on signed-$\dz$ data
        with zero learnable parameters --- below the gen-2
        corrected reference of \SI{125}{\mu m}.  This becomes
        the new gen-3 baseline.
  \item The Gen-2 squared-MSE loss with data-spread normalisation
        is heavy-tail dominated under signed $\dz$ and is unsuitable
        as a model-selection metric in the gen-3 regime
        (factor-1100 inter-split variance vs $<10\,\%$ on
        median residuals).
  \item The Gen-2 correction-net architecture does \emph{not}
        generalise verbatim to signed $\dz$: under the same
        LR schedule it monotonically degrades pure RK4 on
        the body of the distribution.  The required structural
        change is a truncation-error-gated multiplicative
        correction (Fix~N), not a feature addition.
\end{enumerate}

% ============================================================================
\section{Known limitations carried forward}\label{gen3:sec:limits}

\begin{itemize}
  \item \textbf{Material effects --- confounder for G2.}  Baseline
        Gen-3 trains on vacuum propagation; Bethe--Bloch and
        Highland are staged for Gen-3.1.  This is a test-design
        confounder, not just a limitation: a failing G2 at M4 cannot
        be attributed to the extrapolator versus missing $dE/dx$.
        The mitigation adopted here is to run Stage~3 with
        \emph{material effects disabled in the standalone Kalman
        filter at M4}, re-enabling them only at Gen-3.1.  A
        light-weight Bethe--Bloch correction in the generator
        ($dE/dx$ between RK4 sub-steps, makes $qop_f\ne qop_0$ and
        exercises the A5 output head) is kept as a fallback if
        Stage~3 with material-off indicates residual $\delta p/p$
        sensitivity.
  \item \textbf{Single polarity, symmetry assumption.}  MagDown
        only; MagUp would be exact under $qop\to-qop$ only for a
        pure $B_y$ field.  With the real interpolated field map
        ($B_x,B_z$ non-zero in fringes) the symmetry is approximate.
        MagUp training is staged for Gen-3.1; the A6 \texttt{polarity}
        header field ensures a MagUp Allen run loading a MagDown
        binary raises a hard error.
  \item \textbf{Field-map uncertainty.}  Gen-3 treats the
        \texttt{twodip.rtf} field map as ground truth; its measured
        RMS error is ${\sim}1.3\,\%$.  A 1\,\% field error maps
        directly to a $\sim 1\,\%$ $\delta p/p$ bias in G2
        regardless of extrapolator quality; Gate~G2 should be read
        modulo this irreducible floor.
  \item \textbf{Slope-rejection cut.}  The Gen-2
        post-integration rejection $|t|<2.0$ is retained but
        tightened to $|t|<0.5$, comfortably above the LHCb
        geometrical acceptance of ${\sim}0.3$\,rad and the training
        envelope of $\pm 0.4/0.35$.
  \item \textbf{No measurement noise.}  The data generator emits
        exact RK4 trajectories; Allen feeds Kalman-estimated
        states with covariance.  The training-inference mismatch
        on input distribution is mitigated by C3 (wide sampling) but
        not eliminated.
  \item \textbf{Fix~J cost.} Jacobian regulariser is expensive;
        per-epoch training time projected at $2\text{--}3\times$
        baseline.  Acceptable within the long-queue HTCondor
        policy.
\end{itemize}

% ============================================================================
\section{Reproducibility}\label{gen3:sec:repro}

\begin{itemize}
  \item \textbf{Environment.}  \texttt{TE} conda env; PyTorch
        2.9.1+cu128; Python 3.10; identical to Gen-2.
        \texttt{environment.yml} and \texttt{requirements.txt}
        pinned.
  \item \textbf{Code.}  Branch \texttt{gen\_3} of
        \texttt{TrackExtrapolation}.  New modules:
        \texttt{experiments/gen\_3/utils/rk4\_propagator.py}
        (C1-patched), \texttt{utils/magnetic\_field.py} (shared with
        Gen-2), \texttt{models/architectures.py} (Fix H input shape,
        A5 output head), \texttt{models/train.py} (Fix I
        detector-$\sigma$ loss, Fix J Jacobian regulariser).
  \item \textbf{Configs.}  \texttt{experiments/gen\_3/configs/}
        with one JSON per candidate; each declares
        \texttt{applied\_fixes=[C1,C2,C3,A5,H,I,\dots]} and
        \texttt{supersedes=\{gen2\_experiment\_id\}}.
  \item \textbf{MLflow.} Experiment
        \texttt{gen\_3\_track\_extrapolation}; run tags include
        \texttt{applied\_fixes}, \texttt{supersedes},
        \texttt{loader\_version}, \texttt{dataset\_sha256},
        \texttt{generator\_git\_hash}, \texttt{polarity}.
        Metrics (floats, not booleans) include
        \texttt{stage1\_velo\_mean\_dx},
        \texttt{stage1\_ut\_mean\_dx},
        \texttt{stage1\_scifi\_rho\_dx\_qop},
        \texttt{stage1\_backward\_pass\_ratio},
        \texttt{stage2\_jac\_median\_rel\_err},
        \texttt{stage2\_jac\_p99\_rel\_err},
        \texttt{smoketest\_g1\_chi2ndof},
        \texttt{smoketest\_g2\_dpp\_bias},
        \texttt{smoketest\_g3a\_pull\_std},
        \texttt{smoketest\_g3b\_pull\_mean},
        \texttt{throughput\_per\_track\_ns}.  Two passing candidates
        must be rank-orderable from MLflow alone.
  \item \textbf{Data.}  \texttt{train\_50M\_gen3.npz} +
        \texttt{\dots.meta.json}.  SHA-256 logged in MLflow and
        written into every exported Allen binary (A6).
  \item \textbf{Allen binaries.}  Written to
        \texttt{Allen/ML\_research/models/gen3/} with the filename
        pattern \texttt{\{candidate\}\_\{applied\_fixes\}\_\{git\_hash\}.bin}.
\end{itemize}

% ============================================================================
\section{Summary}\label{gen3:sec:summary}

Gen-3 is a \emph{deployment-first} iteration.  The protocol does not
introduce new model families; instead it takes the best Gen-2
candidate per family, patches the five blocking issues identified in
the Gen-2 §5 Allen-readiness audit (C1, C2, C3, A2, A3), adds four
theoretical improvements (Fixes G--J) derived from the Gen-2 §4.5.1
SciFi failure diagnostics, and runs each candidate through a
three-stage validation funnel terminating in the Allen standalone
smoke test.  The primary success metric is binary: at least one
candidate passes G1+G2+G3 simultaneously at milestone M4.

\begin{itemize}
\item \textbf{Primary candidate:}  \texttt{nrk4\_tiny\_1step} with
      Fixes C1+C2+C3+A5+H+I; first to attempt the smoke test
      because it is the only Gen-2 model that satisfies A2 and A3
      simultaneously.
\item \textbf{Accuracy ceiling:}  \texttt{nrk4\_small\_1step} with
      the same fixes plus float16 export.
\item \textbf{Physics-informed rescue:}  \texttt{pinn\_v2\_lam0.1}
      with Fix~G; if the $\rho(\Delta x,qop)$ criterion is not met
      after Fix~G, PINN\_v2 is retired from the deployment path.
\item \textbf{Reference only:}  \texttt{mlp\_medium} with Fix~K
      ablation; not a deployment target.
\end{itemize}

A Gen-3.1 follow-up is staged to address material effects
(Bethe--Bloch, Highland) and mixed magnet polarity once the vacuum
deployment is proven.

\paragraph{Reviewer pass-through.}  This protocol incorporates
blocking issues B1--B4 (C1 coefficient factor-$10^3$ error, Allen
loader V3 contract, Jacobian dimensionality, component-dependent
detector $\sigma$) and advisable items N1--N15 from the
\texttt{experiment-reviewer} audit (2026-04-23); all pre-M0 gates
have been hardened to numerically enforceable thresholds.

% ============================================================================
\begin{thebibliography}{9}

\bibitem{gen1report}
G.~Scriven, \emph{Gen-1 neural track extrapolation for Allen HLT1},
internal note, 2026.

\bibitem{gen2report}
G.~Scriven, \emph{Gen-2 neural track extrapolation on the
\texttt{train\_50M\_dz25} dataset}, internal note,
\texttt{gen2\_new\_data\_results.tex}, 2026.

\bibitem{gen2protocol}
G.~Scriven, \emph{Gen-2 experimental protocol}, internal note,
\texttt{gen2\_experimental\_protocol.tex}, 2026.

\bibitem{AllenREADME}
LHCb Allen team, \emph{ML\_research: machine-learning extrapolator
integration in Allen HLT1}, \texttt{Allen/ML\_research/README.md},
2026.

\bibitem{AllenDeepDive}
G.~Scriven, \emph{Gen-2 true-data Allen deep dive},
\texttt{Allen/ML\_research/gen2\_true\_data\_deep\_dive.md}, 2026.

\end{thebibliography}

% ============================================================================
\section{Audit addendum (2026-05-06): PINN\_v2 / NeuralRK4 code review and
         Allen deployment kernel}\label{gen3:sec:audit-addendum}

A focused source-level audit of
\texttt{models/architectures.py} and the Allen
\texttt{ParKalman/include/} headers, performed on 2026-05-06,
adds three architectural findings that supplement F1--F3 of
\S\ref{gen3:sec:m1-v1-outcome} and crystallises the deployment
kernel design.

\paragraph{PINN\_v2 magnet-kick term verified correct.}  Direct
comparison of the cross-product expressions in the gen-3 PDE
residual against (i) the gen-3 RK4 propagator and (ii)
Allen's \texttt{ExtrapolatorCommon::derivative} confirms that the
gen-2 \S4.5.1 $\rho(\Delta x, qop) = +0.92$ bias does \emph{not}
arise from a sign or factor error in the residual.  It originates
from three structural defects:
\begin{enumerate}\setlength{\itemsep}{0pt}
  \item Dimensionally inhomogeneous PDE weighting
        (\texttt{pde\_scale = output\_std[:4] / |dz|}) mixes
        mm-scale $(x,y)$ with dimensionless $(t_x,t_y)$ entries
        and produces an asymmetric $|\dz|$-weighting that
        dominates the loss in the high-$|qop|$ tail.
  \item The IC envelope guarantees
        $\mathrm{forward\_at\_z}(z_\mathrm{frac}=0)\equiv$
        initial state by construction; the IC loss has no
        gradient and \texttt{lambda\_ic} is operationally inert.
  \item The 4-output head writes independently into
        $(x,y,t_x,t_y)$, so $\mathrm{d}/\mathrm{d}z\,(x_\mathrm{out})\neq t_{x,\mathrm{out}}$ even
        at $z_\mathrm{frac}=1$.  The network can satisfy the
        data loss while violating its own ODE.
\end{enumerate}
PINN\_v2 is therefore retired from the M1.5 deployment path; Fix~G
alone does not stabilise it.  A closed-form unit test on a
uniform $B_y$ helix (sagitta $1.499{\times}10^{-2}$~mm at
$\dz=1$~m, $q/p=1/(10\,\mathrm{GeV})$) is the gate before any PINN
retrain.

\paragraph{NeuralRK4 correction applies inside \texttt{\_rhs} at
every RK4 stage; \emph{not} truncation-gated.}  The current
implementation evaluates the same correction MLP at $k_1,k_2,k_3,k_4$
and rescales by $\exp(s_\mathrm{corr})\cdot
\mathrm{output\_std}/|\dz|_\mathrm{clamp}$.  The scaling is
unbounded: a $\sim 1.5$~nat upward drift in $s_\mathrm{corr}$
drives the per-stage correction to $\mathcal{O}(\mathrm{output\_std})\sim\SI{10}{mm}$,
quantitatively explaining the F3 \SI{7.1}{mm} blow-up.  Fix~N is
implemented by moving the correction \emph{outside} \texttt{\_rhs},
applying it once per RK4 step on the state increment, and gating
it by an FSAL-style estimator
$\sigma=\tanh\!\bigl(\alpha\,\|k_4-k_3\|\,|h|\bigr)$.  Of the three
estimators considered (FSAL, closed-form $|\dz|^5$, Richardson
half-step), FSAL is recommended: zero extra field calls, tracks
the true local truncation error rather than the structural
$|\dz|^5$ proxy, and naturally amplifies near the magnet centre.

\paragraph{Allen deployment kernel.}
\begin{itemize}\setlength{\itemsep}{0pt}
  \item No GPU MLP/NN extrapolator currently exists in Allen;
        \texttt{ML\_research/standalone/MLPExtrapolator.cuh} is
        host-only.  The new kernel
        \texttt{NeuralRK4Extrapolator.cuh} is to live alongside
        \texttt{RungeKuttaExtrapolator.cuh}.
  \item Weights in CUDA \texttt{\_\_constant\_\_} memory (gate~A1).
        Concrete budgets:
        \texttt{nrk4\_tiny\_1step}\ $=$\ 20\,084~B
        ($\SI{31}{\percent}$ of the \SI{64}{kB} A3 budget);
        \texttt{nrk4\_small\_1step} fp32 $=$ 72\,824~B
        (violates A3 by $\SI{14}{\percent}$); fp16 weights
        with fp32 accumulators $=$ 36\,512~B (passes A3 with
        $\SI{43}{\percent}$ headroom).
  \item Field grid \emph{not} shipped with the model:
        \texttt{Magfield::fieldVectorLinearInterpolation(float3)}
        is the on-device API; Allen already carries the
        $1{\times}3{\times}146{\times}81{\times}81$~fp32 grid
        in texture memory.  This closes F1: the \SI{11}{MB} bulk
        of the gen-3 checkpoints is the
        \texttt{field.B\_grid} buffer, which is stripped at
        export.
  \item Latency on A5000: 4 field lookups $+\,\sim\!5\,000$
        fp32 MACs $\approx \SI{0.3}{\mu s}$ per track,
        comfortably inside the single-digit-$\mu$s envelope.
  \item Jacobian (gate~A7): hand-unrolled forward-mode AD
        recommended.  Allen's
        \texttt{RungeKuttaNystromExtrapolator::incrementJacobian}
        already implements the closed-form RK4 path; we extend it
        with a 5-wide tangent through each $\mathrm{Linear}+\tanh$
        of the correction MLP (5$\times$ MLP cost
        $\approx\SI{150}{ns}$).  Finite-difference (5$\times$
        cost; pull-width artefacts) and \texttt{torch.func.jvp}
        (libtorch on-device) are rejected.
\end{itemize}

\paragraph{V3 loader binary contract.}  Little-endian, 32~B fixed
header (\texttt{magic = 0x4E524B34}, \texttt{version = 3},
\texttt{model\_kind} $\in\{\mathrm{MLP},\mathrm{PINN\_v2},\mathrm{NeuralRK4},\mathrm{NeuralRK4{+}FixN}\}$,
$\mathrm{input\_dim}=7$, $\mathrm{output\_dim}=5$,
\texttt{weight\_dtype}, \texttt{activation},
$n_\mathrm{layers}$), per-layer records, normalisation buffers,
NRK4 extras (\texttt{log\_corr\_scale},
$\alpha_\mathrm{FixN}$, $n_\mathrm{rk\_steps}$), SHA-256 of the
training data, git hash, CRC32 of preceding bytes.  V2 files
without magic remain readable for back-compatibility.

\paragraph{Smoke-test bridge.}  A bitwise-identical
Python$\leftrightarrow$C++ comparison on a fixed 100-track set
gates any Allen merge:
$\|y_\mathrm{py}-y_\mathrm{cpp}\|_\infty < \SI{e-6}{mm}$ in
$(x,y)$ and $<10^{-9}$ in $(t_x,t_y)$.

\paragraph{Capacity-only MLP probe completed.}  The planned
4-layer $\times$ 128 MLP retrain (\texttt{mlp\_medium\_v1}) has now
completed on the same 200k-track / 80-epoch CPU protocol.  The
result is materially better than \texttt{mlp\_small\_v1} but still
far from the physics-aware baselines: mean
$|\Delta x|=\SI{5.569}{mm}$, mean $|\Delta y|=\SI{4.318}{mm}$,
RMS position error $=\SI{10.969}{mm}$ on the test split, versus
\SI{18.41}{mm} / \SI{13.86}{mm} / \SI{36.3}{mm} for
\texttt{mlp\_small\_v1}.  This is a $\sim 3.3\times$ gain from
capacity alone, but still $\sim 49\times$ worse than the pure-RK4
baseline (\texttt{nrk4\_tiny\_1step\_v1} at \SI{0.113}{mm}).  The
correct conclusion is that extra MLP capacity helps but does not
address the underlying signed-$\dz$ tail structure; M1.5 effort
should remain concentrated on the NRK4 path.

\paragraph{M1.5 next-action backlog (dependency-ordered).}
\begin{enumerate}\setlength{\itemsep}{0pt}
  \item Replace dimensionally inhomogeneous \texttt{pde\_scale};
        rebuild IC envelope as boundary collocation at
        $z_\mathrm{frac}=1$ (architectures.py L325--373).
  \item Add closed-form \texttt{test\_uniform\_By\_helix} unit
        test before any PINN retrain.
  \item Apply Fix~N (truncation-gated correction outside
        \texttt{\_rhs}) and Fix~O (drop $\mathrm{sign}(\dz)$,
        $z$-score $\log_{10}|\dz|/100$); init
        $s_\mathrm{corr}=-5$ with \texttt{clamp\_max(0)}
        (architectures.py L709--820).
  \item Implement Fix~M ($\log\cosh$ + median-$|\Delta x|$
        early-stop) and Fix~P (separate-LR groups, grad-norm
        guard) in \texttt{models/train.py}.
  \item Write \texttt{scripts/export\_v3.py} (strip
        \texttt{field.B\_grid}; emit V3 binary).
  \item Add \texttt{NeuralRK4Extrapolator.cuh} with
        hand-unrolled forward-mode AD Jacobian.
  \item Patch \texttt{MLPExtrapolator.cuh} to V3
        (magic header, fp32 native, 7-in/5-out).
  \item Smoke-test bridge:
        \texttt{dump\_smoke\_tracks.py},
        \texttt{test\_nrk4\_v3.cpp},
        \texttt{nrk4\_smoke\_compare.ipynb}; merge gate as above.
\end{enumerate}
Items 1--4 are training-side; 5--8 are deployment-side and may
proceed in parallel once the V3 export contract (item 5) is
fixed.

\end{document}
